% References -- AUTO-DRAFTED best-effort entries generated from the in-text
% citations. VERIFY EVERY ENTRY before circulation; entries the drafting
% assistant could not confirm carry [details to verify]. Cases and statutes
% follow the references.
\section*{References}
\begingroup
\setlength{\parindent}{-0.28in}\setlength{\leftskip}{0.28in}\small

Alesina, Alberto, Reza Baqir and William Easterly. 1999. ``Public Goods and Ethnic Divisions.'' \emph{Quarterly Journal of Economics} 114(4): 1243--1284.

Alesina, Alberto, Edward Glaeser and Bruce Sacerdote. 2001. ``Why Doesn't the United States Have a European-Style Welfare State?'' \emph{Brookings Papers on Economic Activity} 2001(2): 187--254.

Alt, James E.\ and Robert C.\ Lowry. 1994. ``Divided Government, Fiscal Institutions, and Budget Deficits: Evidence from the States.'' \emph{American Political Science Review} 88(4): 811--828.

Anzia, Sarah F. 2014. \emph{Timing and Turnout: How Off-Cycle Elections Favor Organized Groups}. Chicago: University of Chicago Press.

Balogh, Brian. 2009. \emph{A Government Out of Sight: The Mystery of National Authority in Nineteenth-Century America}. Cambridge: Cambridge University Press.

Bechard, Allison, Corey Lang and Shanna Pearson-Merkowitz. 2024. [details to verify: ballot design and bond approval.]

Berkman, Michael B.\ and Eric Plutzer. 2005. \emph{Ten Thousand Democracies: Politics and Public Opinion in America's School Districts}. Washington, DC: Georgetown University Press.

Berry, Christopher R. 2009. \emph{Imperfect Union: Representation and Taxation in Multilevel Governments}. Cambridge: Cambridge University Press.

Berry, Christopher R.\ and Jacob E.\ Gersen. 2010. ``The Timing of Elections.'' \emph{University of Chicago Law Review} 77(1): 37--64.

Bohn, Henning and Robert P.\ Inman. 1996. ``Balanced-Budget Rules and Public Deficits: Evidence from the U.S.\ States.'' \emph{Carnegie-Rochester Conference Series on Public Policy} 45: 13--76.

Buchanan, James M.\ and Gordon Tullock. 1962. \emph{The Calculus of Consent}. Ann Arbor: University of Michigan Press.

Bunch, Beverly S. 1991. ``The Effect of Constitutional Debt Limits on State Governments' Use of Public Authorities.'' \emph{Public Choice} 68(1--3): 57--69.

Burns, Nancy. 1994. \emph{The Formation of American Local Governments: Private Values in Public Institutions}. New York: Oxford University Press.

California Legislative Analyst's Office. 2009. [details to verify: report quantifying majority-support school bond failures under the two-thirds requirement and approval rates after Proposition 39.]

Carr, Jered B. 2018. [details to verify: local debt restrictions and evasion.]

Cellini, Stephanie Riegg, Fernando Ferreira and Jesse Rothstein. 2010. ``The Value of School Facility Investments: Evidence from a Dynamic Regression Discontinuity Design.'' \emph{Quarterly Journal of Economics} 125(1): 215--261.

Cole, [details to verify]. 2023. [details to verify: railroad-aid referenda and the constitutional response.]

Dinan, John J. 2006. \emph{The American State Constitutional Tradition}. Lawrence: University Press of Kansas.

Dove, John A. 2014. [details to verify: historical state and local debt restrictions.]

Einhorn, Robin L. 2006. \emph{American Taxation, American Slavery}. Chicago: University of Chicago Press.

Einstein, Katherine Levine, David M.\ Glick and Maxwell Palmer. 2019. \emph{Neighborhood Defenders: Participatory Politics and America's Housing Crisis}. Cambridge: Cambridge University Press.

Fischel, William A. 2001. \emph{The Homevoter Hypothesis: How Home Values Influence Local Government Taxation, School Finance, and Land-Use Policies}. Cambridge, MA: Harvard University Press.

Gerber, Elisabeth R. 1996. ``Legislative Response to the Threat of Popular Initiatives.'' \emph{American Journal of Political Science} 40(1): 99--128.

Gerber, Elisabeth R., Arthur Lupia, Mathew D.\ McCubbins and D.\ Roderick Kiewiet. 2004. [details to verify: \emph{Stealing the Initiative: How State Government Responds to Direct Democracy} (2001), Prentice Hall -- year cited in text as 2004.]

Grinath, Arthur III, John Joseph Wallis and Richard Sylla. 1997. ``Debt, Default, and Revenue Structure: The American State Debt Crisis in the Early 1840s.'' NBER Historical Working Paper 97.

Hacker, Jacob S. 2004. ``Privatizing Risk without Privatizing the Welfare State: The Hidden Politics of Social Policy Retrenchment in the United States.'' \emph{American Political Science Review} 98(2): 243--260.

Hackworth, Jason. 2007. \emph{The Neoliberal City: Governance, Ideology, and Development in American Urbanism}. Ithaca: Cornell University Press.

Hajnal, Zoltan L., Vladimir Kogan and G.\ Agustin Markarian. 2022. ``Who Votes: City Election Timing and Voter Composition.'' \emph{American Political Science Review} 116(1): 374--383.

Hajnal, Zoltan L., Paul G.\ Lewis and Hugh Louch. 2002. \emph{Municipal Elections in California: Turnout, Timing, and Competition}. San Francisco: Public Policy Institute of California.

Hillhouse, Albert M. 1936. \emph{Municipal Bonds: A Century of Experience}. New York: Prentice-Hall.

Hirt, [details to verify]. 2026. [details to verify: officeholder partisanship and municipal capital composition.]

Huang, [details to verify], Munevar and Samuels. 2026. [details to verify: referendum regimes and city debt instruments/yield spreads.]

Jenkins, Destin. 2021. \emph{The Bonds of Inequality: Debt and the Making of the American City}. Chicago: University of Chicago Press.

Kiewiet, D.\ Roderick and Kristin Szakaly. 1996. ``Constitutional Limitations on Borrowing: An Analysis of State Bonded Indebtedness.'' \emph{Journal of Law, Economics, and Organization} 12(1): 62--97.

Kogan, Vladimir, St\'ephane Lavertu and Zachary Peskowitz. 2018. ``Election Timing, Electorate Composition, and Policy Outcomes: Evidence from School Districts.'' \emph{American Journal of Political Science} 62(3): 637--651.

Logan, John R.\ and Harvey L.\ Molotch. 1987. \emph{Urban Fortunes: The Political Economy of Place}. Berkeley: University of California Press.

Lupia, Arthur and John G.\ Matsusaka. 2004. ``Direct Democracy: New Approaches to Old Questions.'' \emph{Annual Review of Political Science} 7: 463--482.

Mahoney, James and Kathleen Thelen, eds. 2010. \emph{Explaining Institutional Change: Ambiguity, Agency, and Power}. Cambridge: Cambridge University Press.

Matsusaka, John G. 2004. \emph{For the Many or the Few: The Initiative, Public Policy, and American Democracy}. Chicago: University of Chicago Press.

McGann, Anthony J. 2004. ``The Tyranny of the Supermajority: How Majority Rule Protects Minorities.'' \emph{Journal of Theoretical Politics} 16(1): 53--77.

Mettler, Suzanne. 2011. \emph{The Submerged State: How Invisible Government Policies Undermine American Democracy}. Chicago: University of Chicago Press.

Milesi-Ferretti, Gian Maria. 2004. ``Good, Bad or Ugly? On the Effects of Fiscal Rules with Creative Accounting.'' \emph{Journal of Public Economics} 88(1--2): 377--394.

Morris, [details to verify]. 1958. [details to verify: Yale Law Journal article on public building authorities; text elsewhere dates the piece 1962 -- reconcile.]

Mullins, Daniel R.\ and Bruce A.\ Wallin. 2004. ``Tax and Expenditure Limitations: Introduction and Overview.'' \emph{Public Budgeting \& Finance} 24(4): 2--15.

Oliver, J.\ Eric. 2001. \emph{Democracy in Suburbia}. Princeton: Princeton University Press.

Orren, Karen and Stephen Skowronek. 2004. \emph{The Search for American Political Development}. Cambridge: Cambridge University Press.

Peterson, Paul E. 1981. \emph{City Limits}. Chicago: University of Chicago Press.

Pierson, Paul. 2004. \emph{Politics in Time: History, Institutions, and Social Analysis}. Princeton: Princeton University Press.

Poterba, James M. 1994. ``State Responses to Fiscal Crises: The Effects of Budgetary Institutions and Politics.'' \emph{Journal of Political Economy} 102(4): 799--821.

Primo, David M. 2007. \emph{Rules and Restraint: Government Spending and the Design of Institutions}. Chicago: University of Chicago Press.

Rivera, [details to verify]. 2016. [details to verify: school district lease and certificate financing.]

Romer, Thomas and Howard Rosenthal. 1978. ``Political Resource Allocation, Controlled Agendas, and the Status Quo.'' \emph{Public Choice} 33(4): 27--43.

Romer, Thomas and Howard Rosenthal. 1982. ``Median Voters or Budget Maximizers: Evidence from School Expenditure Referenda.'' \emph{Economic Inquiry} 20(4): 556--578.

Sbragia, Alberta M. 1996. \emph{Debt Wish: Entrepreneurial Cities, U.S.\ Federalism, and Economic Development}. Pittsburgh: University of Pittsburgh Press.

Schwartzberg, Melissa. 2014. \emph{Counting the Many: The Origins and Limits of Supermajority Rule}. Cambridge: Cambridge University Press.

Stone, Clarence N. 1989. \emph{Regime Politics: Governing Atlanta, 1946--1988}. Lawrence: University Press of Kansas.

Streeck, Wolfgang and Kathleen Thelen, eds. 2005. \emph{Beyond Continuity: Institutional Change in Advanced Political Economies}. Oxford: Oxford University Press.

Tarr, G.\ Alan. 1998. \emph{Understanding State Constitutions}. Princeton: Princeton University Press.

Tiebout, Charles M. 1956. ``A Pure Theory of Local Expenditures.'' \emph{Journal of Political Economy} 64(5): 416--424.

Trounstine, Jessica. 2018. \emph{Segregation by Design: Local Politics and Inequality in American Cities}. Cambridge: Cambridge University Press.

Tsebelis, George. 2002. \emph{Veto Players: How Political Institutions Work}. Princeton: Princeton University Press.

Wallis, John Joseph. 2005. ``Constitutions, Corporations, and Corruption: American States and Constitutional Change, 1842 to 1852.'' \emph{Journal of Economic History} 65(1): 211--256.

Yu, [details to verify], Chen and Robbins. 2022. [details to verify: bond referenda and secondary-market yields of school district debt.]

Zackin, Emily. 2013. \emph{Looking for Rights in All the Wrong Places: Why State Constitutions Contain America's Positive Rights}. Princeton: Princeton University Press.

\medskip
\noindent\textbf{Cases and statutes cited.}\ \emph{Gelpcke v.\ Dubuque}, 68 U.S.\ (1 Wall.) 175 (1863); \emph{Cipriano v.\ City of Houma}, 395 U.S.\ 701 (1969); \emph{City of Phoenix v.\ Kolodziejski}, 399 U.S.\ 204 (1970); \emph{Salyer Land Co.\ v.\ Tulare Lake Basin Water Storage District}, 410 U.S.\ 719 (1973); \emph{Ball v.\ James}, 451 U.S.\ 355 (1981); \emph{Rose v.\ Council for Better Education}, 790 S.W.2d 186 (Ky.\ 1989); Iowa Code \S\S 384.24--384.26.

\endgroup
